Write \(u_{y,z}=\lambda_{(y,0,z)}\) and \(v_{y,z}=\lambda_{(y,1,z)}\), and order each arm as \((u_{0,z},u_{1,z},v_{0,z},v_{1,z})\). For \(n=2\), \(a=1\), and \(s=(2,3)\), the definition of \(w(a,s)\) gives, before treatment swapping and sign reversal, the three blocks \((-1,-1,-1,0)\), \((1,0,0,0)\), and \((0,0,0,0)\). Therefore \(\bar w(1,(2,3))\) has blocks \((1,0,1,1)\), \((0,0,-1,0)\), and \((0,0,0,0)\). Its two coordinates fixed by the gauge already vanish, so
\[
\rho=\Gamma(\bar w(1,(2,3)))=\bigl((1,0,1,1),(0,0,-1,0),(0,0,0,0)\bigr).
\]
Define the second gauge-fixed vector
\[
\xi=\Gamma(\bar w(1,(3,2)))=\bigl((1,0,1,1),(0,0,0,0),(0,0,-1,0)\bigr).
\]
The same calculation from the definition of \(w\), with arms two and three interchanged, proves the second equality. In particular, both vectors obey \(u_{0,2}=u_{0,3}=0\).

For \(\alpha\in\{\rho,\xi\}\), a treatment word \(\delta\in\{0,1\}^3\), and an outcome-response pair \(\eta=(a,b)\in\{0,1\}^2\), define the dual slack
\[
s^{\alpha}_{\delta,ab}=\sum_{z=1}^3\alpha_{(\eta(\delta(z)),\delta(z),z)}-(b-a).
\]
This is precisely \((A_{2,3}^{\mathsf T}\alpha)_{(\delta,\eta)}-(c_{2,3})_{(\delta,\eta)}\) by the incidence and ATE-loading definitions. Direct substitution gives all thirty-two slacks for each vector:
\[
\begin{array}{c|rrrr}
\delta\backslash\eta&00&01&10&11\\ \hline
000&1&0&1&0\\
001&1&0&1&0\\
010&0&0&0&0\\
011&0&0&0&0\\
100&1&0&2&1\\
101&1&0&2&1\\
110&0&0&1&1\\
111&0&0&1&1
\end{array}
\qquad(\alpha=\rho),
\]
\[
\begin{array}{c|rrrr}
\delta\backslash\eta&00&01&10&11\\ \hline
000&1&0&1&0\\
001&0&0&0&0\\
010&1&0&1&0\\
011&0&0&0&0\\
100&1&0&2&1\\
101&0&0&1&1\\
110&1&0&2&1\\
111&0&0&1&1
\end{array}
\qquad(\alpha=\xi).
\]
Every displayed slack is nonnegative. Hence \(\rho,\xi\in\mathcal H^0_{2,3}\), and the zero entries in the two tables give their complete signatures.

Use the lexicographic type order prescribed by the handle: first the three-bit word \(\delta\), then the pair \((a,b)\). Scan the zero entries in that order and retain a row whenever its normal increases the rank in the gauge coordinates. The cumulative ranks for \(\rho\), in the order of its complete signature, are
\[
(1,2,3,4,5,6,7,7,8,8,8,8,9,9,10,10,10,10),
\]
and those for \(\xi\) are
\[
(1,2,3,4,5,5,6,7,8,8,8,8,9,10,10,10,10,10).
\]
Consequently the greedy scans return
\[
\begin{aligned}
B_{\rho}=(&000/01,000/11,001/01,001/11,010/00,010/01,\\
          &010/10,011/00,100/01,110/00),\\
B_{\xi}=(&000/01,000/11,001/00,001/01,001/10,010/01,\\
         &010/11,011/00,100/01,101/00),
\end{aligned}
\]
where \(\delta/ab\) abbreviates \((\delta,(a,b))\). For completeness, the rank calculation is certified over the integers as follows. In the reduced coordinate order
\[
x=(u_{0,1},u_{1,1},v_{0,1},v_{1,1},u_{1,2},v_{0,2},v_{1,2},u_{1,3},v_{0,3},v_{1,3}),
\]
the row-normal matrices of the two displayed bases are
\[
R_{\rho}=\begin{pmatrix}
1&0&0&0&0&0&0&0&0&0\\
0&1&0&0&1&0&0&1&0&0\\
1&0&0&0&0&0&0&0&0&1\\
0&1&0&0&1&0&0&0&0&1\\
1&0&0&0&0&1&0&0&0&0\\
1&0&0&0&0&0&1&0&0&0\\
0&1&0&0&0&1&0&1&0&0\\
1&0&0&0&0&1&0&0&1&0\\
0&0&0&1&0&0&0&0&0&0\\
0&0&1&0&0&1&0&0&0&0
\end{pmatrix},
\qquad \det(R_{\rho})=-1,
\]
\[
R_{\xi}=\begin{pmatrix}
1&0&0&0&0&0&0&0&0&0\\
0&1&0&0&1&0&0&1&0&0\\
1&0&0&0&0&0&0&0&1&0\\
1&0&0&0&0&0&0&0&0&1\\
0&1&0&0&1&0&0&0&1&0\\
1&0&0&0&0&0&1&0&0&0\\
0&1&0&0&0&0&1&1&0&0\\
1&0&0&0&0&1&0&0&1&0\\
0&0&0&1&0&0&0&0&0&0\\
0&0&1&0&0&0&0&0&1&0
\end{pmatrix},
\qquad \det(R_{\xi})=1.
\]
Thus each vector is isolated in the ten-dimensional gauge section by ten tight type equations. Since each is feasible, each is a vertex of \(\mathcal H^0_{2,3}\); in particular, \(\rho,\xi\in\mathcal V_{2,3}\). The displayed cumulative ranks also show that every earlier tight row omitted by the scan is dependent on the rows already retained. At any scan position, a dependent row cannot enlarge the current independent prefix, whereas an independent row must be included in the lexicographically least full-rank extension because every possible replacement occurs later. Induction over the ordered signature therefore proves that \(B_{\rho}\) and \(B_{\xi}\) are exactly the canonical bases specified by the reverse-search handle.

The first two entries of \(B_{\rho}\) and \(B_{\xi}\) agree. Their third entries are respectively \(001/01\) and \(001/00\), and \(001/00\prec001/01\). Hence
\[
B_{\xi}\prec B_{\rho}.
\]
Every parent edge prescribed by the handle strictly decreases the canonical basis. It follows inductively that every basis on a parent chain starting at \(\xi\), after its initial basis, is strictly smaller than \(B_{\xi}\), and therefore strictly smaller than \(B_{\rho}\). No such chain can reach \(\rho\). Thus the required root-reachability property fails, so the specified decreasing-basis rule does not form a reverse-search tree rooted at \(\rho\) and cannot establish the claimed polynomial-delay enumeration. The argument concerns only this parent rule; it gives no obstruction to a different polynomial-delay algorithm.